\documentclass[11pt,a4paper]{article}
% XeLaTeX-friendly preamble: use fontspec for Unicode fonts
\usepackage{fontspec}
\setmainfont{TeX Gyre Termes}
\usepackage[brazil]{babel}
\usepackage{geometry}
\geometry{margin=2.5cm}
\usepackage{amsmath,amssymb}
\usepackage{enumitem}
\usepackage{hyperref}
\usepackage{longtable}
\usepackage{graphicx}
\usepackage{float}
\usepackage{booktabs}
\usepackage{microtype}
\usepackage{fancyvrb}
% Use a Unicode monospaced font for verbatim/code blocks (improves glyph coverage)
\setmonofont{DejaVu Sans Mono}[Scale=MatchLowercase]
\fvset{fontsize=\small,formatcom=\ttfamily}
\usepackage{comment}
\title{Revisão para a prova — Engenharia de Reservatórios I}
\author{Compilado}
\date{Gerado em 2026-04-22}
\begin{document}
\maketitle
\tableofcontents
\clearpage

% ---- file: $cap1_{transcription}$.md
\section{ENGENHARIA DE RESERVATÓRIOS I}
 \textbf{3º Ano – 2º Semestre — 2025–2026}

 \textbf{Docente:} Prof. Dr. Geraldo A. R. Ramos

 ---

\subsection{Sumário}
\begin{itemize}
\item 1. Sistema de Produção
\item 2. Engenharia de Reservatórios: funções e objetivos
\item 3. História da Engenharia de Reservatórios
\item 4. Tipos de reservatórios (diagramas de fases)
\item 5. Mecanismos de produção
\item 6. Consolidação da aula (exercícios)
\item 7. Tarefas
\item 8. Trabalho investigativo
\item 9. Bibliografia
\end{itemize}

 ---

\subsection{Objetivo geral}
 Analisar os princípios fundamentais da Engenharia de Reservatórios, abrangendo sistemas de produção, funções e objetivos da disciplina, evolução histórica, classificação de reservatórios segundo diagramas de fases e mecanismos de produção, visando a capacitação para integração de informações e apoio à tomada de decisão técnica em campos petrolíferos.

\subsubsection{Objetivos específicos}
\begin{itemize}
\item Descrever os componentes e a estrutura de um sistema de produção de hidrocarbonetos.
\item Explicar as principais funções e objetivos da Engenharia de Reservatórios na exploração e desenvolvimento de campos petrolíferos.
\item Relatar a evolução histórica da disciplina, destacando marcos tecnológicos e metodológicos.
\item Classificar os tipos de reservatórios com base em diagramas de fases (óleo, gás, condensado).
\item Comparar os principais mecanismos de produção: drive por pressão, expansão de gás, influxo de água e outros processos naturais.
\end{itemize}

\subsection{Sistema de Produção}
 O sistema de produção é o conjunto de equipamentos destinados à coleta, separação, tratamento, estocagem e escoamento dos fluidos produzidos e movimentados no campo de petróleo ou gás natural. A produção ocorre em três etapas principais:

\subsubsection{Etapas}
\begin{itemize}
\item Recuperação: fluxo de fluidos no reservatório até o fundo do poço.
\item Elevação: fluxo do fundo do poço até a superfície (topo do poço), que pode ser natural ou artificial.
\item Escoamento: fluxo do topo do poço até o separador, passando por reguladores de fluxo.
\end{itemize}

\subsubsection{Elementos do sistema petrolífero}
\begin{itemize}
\item \textbf{Armadilha (trapa):} arranjo estrutural que permite a acumulação de hidrocarbonetos (dobras, falhas, fraturas).
\item \textbf{Rocha geradora:} rocha rica em matéria orgânica que gera hidrocarbonetos.
\item \textbf{Rocha selante:} rocha com baixa permeabilidade que impede a migração de hidrocarbonetos.
\item \textbf{Migração:} processo que mobiliza os hidrocarbonetos da zona geradora para o reservatório.
\item \textbf{Sincronismo:} coincidência temporal dos processos geológicos necessários para formar um sistema petrolífero.
\item \textbf{Soterramento e catagênese:} processos térmicos que transformam querogênio em óleo, condensado ou gás.
\end{itemize}

 \emph{Nota:} Querogênio é a fração insolúvel da matéria orgânica; catagênese refere-se ao craqueamento térmico das moléculas maiores em moléculas menores sob aumento de temperatura e pressão.

\subsection{Engenharia de Reservatórios — Definição}
 A Engenharia de Reservatórios é o ramo da Engenharia do Petróleo que estuda o fluxo de fluidos (hidrocarbonetos) no interior das rochas porosas, visando maximizar o fator de recuperação. É uma atividade multidisciplinar que envolve física, química, geologia, engenharia, matemática e ciência da computação.

\subsection{Consolidação da aula — Exercícios}
 Responda as questões abaixo e justifique suas escolhas.

\begin{enumerate}
\item O óleo produzido em reservatórios leves caracteriza-se por:
\begin{itemize}
\item A) Alta viscosidade e baixa mobilidade
\item B) Baixa densidade e alta mobilidade
\item C) Formação de gás livre em excesso
\item D) Saturação irreductível de água igual a zero
\item E) Exclusivamente líquido sem variação de densidade
\end{itemize}
\end{enumerate}

\begin{enumerate}
\item O gás produzido em reservatórios influencia a recuperação de óleo porque:
\begin{itemize}
\item A) Reduz a pressão do reservatório e facilita o fluxo de óleo
\item B) Aumenta a densidade do óleo
\item C) Impede totalmente a produção de água
\item D) É irrelevante para a eficiência de produção
\item E) A saturação de óleo se torna irrelevante
\end{itemize}
\end{enumerate}

\begin{enumerate}
\item Reservatórios com mecanismo de water drive caracterizam-se por:
\begin{itemize}
\item A) Manutenção da pressão por influxo de água do aquífero
\item B) Elevação artificial de óleo via bombeio mecânico
\item C) Expansão de gás dissolvido no óleo
\item D) Produção exclusivamente de gás condensado
\item E) Inexistência de óleo saturado
\end{itemize}
\end{enumerate}

\begin{enumerate}
\item A viscosidade do óleo impacta diretamente:
\begin{itemize}
\item A) A taxa de fluxo pelo reservatório
\item B) A pressão do aquífero
\item C) A densidade da água
\item D) A compressibilidade da rocha
\item E) A saturação irreductível de gás
\end{itemize}
\end{enumerate}

\begin{enumerate}
\item Um mecanismo de produção por expansão de gás livre (gas-cap drive):
\begin{itemize}
\item A) Mantém pressão por expansão de gás acima do óleo
\item B) Depende exclusivamente de água injetada
\item C) Gera aumento da viscosidade do óleo
\item D) Reduz a saturação irreductível de água
\item E) Impede o fluxo de óleo
\end{itemize}
\end{enumerate}

\begin{enumerate}
\item O condensado produzido em reservatórios ocorre quando:
\begin{itemize}
\item A) O líquido condensado se separa do gás à medida que a pressão diminui
\item B) Óleo leve é produzido sem formação de gás
\item C) A água do reservatório se transforma em gás
\item D) A pressão está acima da pressão de saturação do gás
\item E) Não há variação de densidade com a pressão
\end{itemize}
\end{enumerate}

\begin{enumerate}
\item Qual das seguintes características NÃO define o mecanismo solution gas drive?
\begin{itemize}
\item A) Expansão do gás dissolvido no óleo
\item B) Redução gradual da pressão do reservatório
\item C) Formação de gás livre a partir do óleo
\item D) Dependência da pressão inicial para mobilidade
\item E) Produção contínua de óleo
\end{itemize}
\end{enumerate}

\begin{enumerate}
\item Em um reservatório com influxo de água (water influx), a produção de óleo é sustentada porque:
\begin{itemize}
\item A) A água do aquífero desloca o óleo em direção aos poços
\item B) A expansão do gás acima do óleo aumenta a viscosidade
\item C) Poços artificiais elevam a pressão
\item D) O óleo condensado impede a produção de água
\item E) A saturação irreductível de óleo é zero
\end{itemize}
\end{enumerate}

\begin{enumerate}
\item A produção de óleo em um reservatório undersaturated depende principalmente de:
\begin{itemize}
\item A) Pressão inicial do reservatório e mobilidade do óleo
\item B) Expansão do gás acima do óleo
\item C) Injeção de água ou gás
\item D) Formação de condensado
\item E) Elevação artificial exclusivamente
\end{itemize}
\end{enumerate}

\begin{enumerate}
\item O que caracteriza um reservatório com gas-cap drive?
\begin{itemize}
\item A) Gás livre no topo do reservatório que ajuda a manter a pressão
\item B) Água injetada artificialmente para suporte de pressão
\item C) Óleo pesado não fluido
\item D) Produção apenas de condensado líquido
\item E) Reservatório totalmente saturado sem gás
\end{itemize}
\end{enumerate}

\subsection{Tarefas}
\begin{enumerate}
\item Explique a diferença entre Sistema de Produção, Sistema Petrolífero e cadeia produtiva (cadeia de valor) da indústria petrolífera.
\item Faça uma ilustração no envelope de fases com os elementos mencionados no item 2.
\item Investigue as três principais teorias sobre a origem dos hidrocarbonetos:
\begin{itemize}
\item a) Qual a diferença entre as teorias?
\item b) Qual das teorias é mais aceita?
\item Qual é a área de atuação do Engenheiro de Reservatórios? Justifique.
\item Complete: HIDROCARBONETOS + NÃO HIDROCARBONETOS = \rule{6cm}{0.5pt}
\item Fundamente as respostas corretas e incorretas dos exercícios de consolidação.
\item Explique as principais propriedades dos fluidos produzidos em reservatórios e como influenciam a produção.
\item Descreva os mecanismos naturais de produção de um reservatório.
\item Compare os efeitos da expansão de gás e do influxo de água sobre a recuperação de óleo.
\item Discuta a importância de integrar informações sobre fluidos e mecanismos de produção no planejamento de campo.
\item Preencha a tabela sobre tipos de reservatórios, mecanismos de produção, condições iniciais, comportamento com o declínio de pressão e hidrocarbonetos produzidos.
\end{itemize}
\end{enumerate}

\subsection{Trabalho investigativo — História da Engenharia de Reservatórios}
 Os estudantes deverão elaborar um trabalho investigativo sobre a temática “História da Engenharia de Reservatórios”. O trabalho deve:
\begin{itemize}
\item Ser estruturado segundo as normas de um TCC: capa, resumo, sumário, introdução, desenvolvimento, conclusão e referências bibliográficas.
\item Ser produzido e apresentado em LaTeX, garantindo formatação acadêmica adequada.
\item Seguir rigorosamente as normas de citação e referência bibliográfica recomendadas pelo curso.
\end{itemize}

 \textbf{Data de submissão:} 24 de março de 2026.

\subsection{Definições importantes}
\begin{itemize}
\item \textbf{Sistema de Produção:} conjunto de instalações que permite a recuperação, elevação e escoamento dos fluidos produzidos desde o reservatório até instalações de superfície.
\item \textbf{Armadilha (trapa):} estrutura geológica que impede a migração de hidrocarbonetos e permite sua acumulação.
\item \textbf{Porosidade} (@@$PLACEHOLDER_{0}$@@): fração volumétrica de vazios na rocha, expressa como percentagem ou fração (por exemplo, 0.18 = 18\%).
\item \textbf{Permeabilidade} (@@$PLACEHOLDER_{1}$@@): medida da facilidade com que um fluido atravessa a rocha (unidade: Darcy ou m@@$PLACEHOLDER_{2}$@@).
\item \textbf{Fator de volume de formação} (@@$PLACEHOLDER_{3}$@@): relação entre o volume de óleo em condições de reservatório e o volume correspondente à superfície (Bo = @@$PLACEHOLDER_{4}$@@).
\item \textbf{Razão de solubilidade} (@@$PLACEHOLDER_{5}$@@): quantidade de gás dissolvido por unidade de óleo (usada em m@@$PLACEHOLDER_{6}$@@/m@@$PLACEHOLDER_{7}$@@ ou scf/STB).
\end{itemize}

\subsection{Exemplo prático — cálculo rápido de OOIP (fórmula de campo)}
Considere: Área = 100 acres; espessura neta h = 20 ft; porosidade @@$PLACEHOLDER_{0}$@@; saturação irreducível de água @@$PLACEHOLDER_{1}$@@; @@$PLACEHOLDER_{2}$@@ bbl/STB.

Usando a fórmula prática:
@@$PLACEHOLDER_{0}$@@
Substituindo valores:
@@$PLACEHOLDER_{1}$@@
Interpretação: aproximadamente 1.45 milhões de barris em lugar antes da produção.

\subsection{Glossário de símbolos (rápido)}
\begin{itemize}
\item @@$PLACEHOLDER_{8}$@@: área (acres ou m@@$PLACEHOLDER_{9}$@@)
\item @@$PLACEHOLDER_{10}$@@: espessura (ft ou m)
\item @@$PLACEHOLDER_{11}$@@: porosidade (fração)
\item @@$PLACEHOLDER_{12}$@@, @@$PLACEHOLDER_{13}$@@: saturação de água, saturação irreducível de água
\item @@$PLACEHOLDER_{14}$@@, @@$PLACEHOLDER_{15}$@@: fatores de volume de formação (óleo, gás)
\item @@$PLACEHOLDER_{16}$@@: óleo originalmente em lugar (OOIP)
\item @@$PLACEHOLDER_{17}$@@, @@$PLACEHOLDER_{18}$@@: produção acumulada de óleo e gás
\item @@$PLACEHOLDER_{19}$@@: razão gás/óleo dissolvido
\end{itemize}

\subsection{Dicas rápidas de estudo}
\begin{itemize}
\item Extraia os símbolos das equações e mantenha uma tabela pessoal de unidades (campo vs SI).
\item Para exercícios volumétricos, verifique unidades antes de aplicar fatores práticos (7758, 43560, etc.).
\item Use pequenos exemplos numéricos para validar o entendimento antes de resolver casos maiores.
\end{itemize}


Muito obrigada!

\subsection{Perguntas modelo — Capítulo 1}
\begin{enumerate}
\item Verdadeiro ou falso: indique V (verdadeiro) ou F (falso) e justifique em até duas linhas.
\begin{enumerate}[label=\alph*)]
\item Em reservatórios leves, a viscosidade do óleo costuma ser menor que em óleos pesados. \hfill \textbf{(V/F)}\\
\textbf{Justifique:} \vspace{1\baselineskip}
\item A presença de gás dissolvido no óleo sempre aumenta a mobilidade do fluido. \hfill \textbf{(V/F)}\\
\textbf{Justifique:} \vspace{1\baselineskip}
\item Uma drive de água (water drive) mantém a pressão do reservatório por influxo de água. \hfill \textbf{(V/F)}\\
\textbf{Justifique:} \vspace{1\baselineskip}
\item A saturação relativa da água é igual a 1 menos a saturação do óleo. \hfill \textbf{(V/F)}\\
\textbf{Justifique:} \vspace{1\baselineskip}
\end{enumerate}
\item Questões curtas (responda em 3–5 linhas):
\begin{enumerate}[label=\alph*)]
\item Descreva brevemente o papel de uma armadilha geológica na acumulação de hidrocarbonetos.
\item Explique a diferença conceitual entre Sistema de Produção e Sistema Petrolífero.
\end{enumerate}
\end{enumerate}


% ---- file: $cap2_{transcription}$.md


\section{ENGENHARIA DE RESERVATÓRIOS I}
\textbf{3º Ano – 2º Semestre — 2025–2026}

\textbf{Docente:} Prof. Dr. Geraldo A. R. Ramos


\subsection{Sumário}
\begin{itemize}
\item 2.1 Introdução
\item 2.2 Misturas e soluções
\item 2.3 Propriedades básicas dos fluidos
\item 2.3.1 Propriedades de óleo
\item 2.3.2 Propriedades dos gases
\item 2.3.3 Propriedades de misturas de hidrocarbonetos
\item 2.4 Consolidação da aula
\item 2.5 Tarefas
\item Bibliografia
\end{itemize}


\subsection{Objetivo geral}
Analisar as propriedades dos fluidos de reservatório por meio da caracterização de misturas e soluções, aplicação de equações de estado, interpretação de dados PVT e utilização de correlações empíricas, visando prever o comportamento dos fluidos e apoiar a tomada de decisão na engenharia de reservatórios.

\subsubsection{Objetivos específicos}
\begin{itemize}
\item Identificar os tipos de misturas e soluções presentes nos fluidos.
\item Descrever propriedades físicas e termodinâmicas relevantes (densidade, viscosidade, compressibilidade, fator de volume, etc.).
\item Aplicar equações de estado e correlações empíricas para estimativas de propriedades.
\item Calcular o fator volume de formação (@@$PLACEHOLDER_{20}$@@), a razão de solubilidade (@@$PLACEHOLDER_{21}$@@) e interpretar resultados PVT.
\item Avaliar a influência da viscosidade e do grau API no escoamento.
\end{itemize}

\subsection{2.1 Introdução}
Um fluido é uma substância que se deforma continuamente (escoa) sob ação de uma força tangencial. Para o engenheiro de reservatórios, conhecer as propriedades físico‑químicas e o comportamento de fase é essencial para estimar volumes in situ, planejar produção e definir estratégias de recuperação.

\subsubsection{Amostragem e análises}
\begin{itemize}
\item Amostragem em superfície (separadores e frascos de amostra).
\item Amostragem em fundo de poço (downhole sampling).
\item Amostragem em teste de formação (formation testing).
\end{itemize}

Análises normalmente realizadas:
\begin{itemize}
\item Determinação da composição (frações de hidrocarbonetos C1–Cn).
\item Análises PVT: ensaios CCE (constant composition expansion), DL (differential liberation), flash, etc.
\item Determinação de propriedades tecnológicas: @@$PLACEHOLDER_{22}$@@, @@$PLACEHOLDER_{23}$@@, viscosidade, densidade, Z‑factor.
\item Testes de compatibilidade de fluidos e análise de contaminantes.
\end{itemize}

\subsection{2.2 Misturas e soluções}
\begin{itemize}
\item Mistura: combinação física de componentes (gás + óleo + água) onde as fases podem coexistir.
\item Solução: situação em que um componente está dissolvido (ex.: gás dissolvido no óleo). A razão gás/óleo dissolvido é expressa por @@$PLACEHOLDER_{24}$@@ (m³ de gás por m³ de óleo, ou SCCM/Sm³ conforme unidade).
\item Pontos relevantes: ponto de bolha (bubble point) e ponto de orvalho (dew point) que definem o início da formação de uma fase livre.
\end{itemize}

\subsection{2.3 Propriedades básicas dos fluidos}

\subsubsection{2.3.1 Propriedades de óleo}
\begin{itemize}
\item Densidade e gravidade API: @@$PLACEHOLDER_{25}$@@ (forma prática para converter densidade específica em grau API).
\item Viscosidade: afeta diretamente a mobilidade e a taxa de fluxo através do meio poroso.
\item Fator volume de formação (@@$PLACEHOLDER_{26}$@@): relação entre volume no reservatório e volume à superfície; importante para conversão de volumes de reservatório para superficie.
\item Compressibilidade de óleo: influencia o declínio de pressão e estimativas volumétricas.
\end{itemize}

\subsubsection{2.3.2 Propriedades dos gases}
\begin{itemize}
\item Gases ideais obedecem à equação de estado @@$PLACEHOLDER_{27}$@@ em condições de baixa pressão e alta temperatura.
\item Leis fundamentais: Boyle ( @@$PLACEHOLDER_{28}$@@ ), Charles ( @@$PLACEHOLDER_{29}$@@ ), Avogadro.
\item Lei de Dalton (pressões parciais): a pressão total de uma mistura é a soma das pressões parciais.
\item Gases reais: desvio do comportamento ideal é quantificado pelo fator de compressibilidade @@$PLACEHOLDER_{30}$@@ (tabelas/curvas de correção). Equações de estado cúbicas (Peng‑Robinson, Soave‑Redlich‑Kwong) são usadas em PVT.
\end{itemize}

\subsubsection{2.3.3 Propriedades de misturas de hidrocarbonetos}
\begin{itemize}
\item Comportamento de fases (monofásico vs bifásico) dependente de pressão e temperatura.
\item Razão gás‑óleo (@@$PLACEHOLDER_{31}$@@) e fator de formação @@$PLACEHOLDER_{32}$@@ em sistemas bifásicos.
\item Uso de análises PVT (flash/diferencial) para determinar curvas de @@$PLACEHOLDER_{33}$@@, @@$PLACEHOLDER_{34}$@@, viscosidade vs pressão e temperatura.
\end{itemize}

\subsection{2.4 Consolidação da aula}
Exemplos de perguntas para revisão:
\begin{itemize}
\item O que é @@$PLACEHOLDER_{35}$@@ e por que é importante para o cálculo de volumes produzidos?
\item Como o @@$PLACEHOLDER_{36}$@@ varia com a pressão e o que significa atingir o ponto de bolha?
\item Como a viscosidade do óleo influencia a mobilidade e o fator de recuperação?
\item Quando o gás pode ser tratado como ideal e quando deve‑se usar correções (Z‑factor)?
\end{itemize}

\subsection{2.5 Tarefas}
\begin{enumerate}
\item Desenhar o envelope de fases para uma mistura óleo‑gás e identificar ponto de bolha e ponto de orvalho.
\item Calcular @@$PLACEHOLDER_{37}$@@ e @@$PLACEHOLDER_{38}$@@ em um caso hipotético a partir de dados PVT simplificados (fornecerei os dados se desejar).
\item Fazer um resumo comparativo: gases ideais vs gases reais e apresentar quando usar equações de estado cúbicas.
\item Ler um ensaio PVT (CCE + DL) e resumir os principais resultados e curvas obtidas.
\end{enumerate}

\subsection{Definições importantes}
\begin{itemize}
\item \textbf{Fator de volume de formação (@@$PLACEHOLDER_{39}$@@):} razão entre o volume de óleo em condições de reservatório e o volume correspondente à superfície (bbl/STB ou m@@$PLACEHOLDER_{40}$@@/m@@$PLACEHOLDER_{41}$@@). Fórmula: @@$PLACEHOLDER_{42}$@@.
\item \textbf{Razão de solubilidade (@@$PLACEHOLDER_{43}$@@):} quantidade de gás dissolvido por unidade de óleo (scf/STB ou m$^3$/m$^3$). No ponto de bolha, @@$PLACEHOLDER_{44}$@@ corresponde ao máximo gás dissolvido antes do aparecimento de gás livre.
\item \textbf{Fator de compressibilidade do gás (@@$PLACEHOLDER_{45}$@@):} correção adimensional que quantifica o desvio do comportamento ideal do gás: @@$PLACEHOLDER_{46}$@@.
\item \textbf{Ponto de bolha / ponto de orvalho:} pressão (a uma dada temperatura) em que surge a primeira bolha de gás no líquido ou a primeira gota de líquido no gás.
\end{itemize}

\subsection{Exemplo prático — cálculo de @@$PLACEHOLDER_{47}$@@ e @@$PLACEHOLDER_{48}$@@}
Considere um ensaio PVT simplificado com resultados experimentais por 1 STB amostrado:
\begin{itemize}
\item Volume de óleo medido em condições de reservatório: @@$PLACEHOLDER_{49}$@@;
\item Volume correspondente à superfície: @@$PLACEHOLDER_{50}$@@;
\item Gás libertado no ensaio: @@$PLACEHOLDER_{51}$@@.
\end{itemize}

Cálculo de @@$PLACEHOLDER_{52}$@@:
@@$PLACEHOLDER_{2}$@@
Interpretação: 1 STB de óleo de superfície corresponde a 1.20 bbl no reservatório.

Conversão de @@$PLACEHOLDER_{53}$@@ para unidades SI (opcional):
1 scf = 0.0283168~m@@$PLACEHOLDER_{54}$@@; 1 STB = 0.1589873~m@@$PLACEHOLDER_{55}$@@.
@@$PLACEHOLDER_{3}$@@

Observação: prefira as unidades pedidas no enunciado (scf/STB é muito comum). Converta unidades apenas quando necessário e mantenha consistência.

\subsection{Notas práticas e correlações}
\begin{itemize}
\item Quando não existirem dados experimentais, use correlações empíricas (p.ex. Standing, Vazquez‑Beggs) ou equações de estado cúbicas (Peng‑Robinson, SRK) para estimar @@$PLACEHOLDER_{56}$@@, @@$PLACEHOLDER_{57}$@@ e viscosidade.
\item Para gás, determine @@$PLACEHOLDER_{58}$@@ via charts ou equações para calcular volumes em condições padrão e avaliar expansibilidade.
\end{itemize}

\subsection{Glossário de símbolos (cap.2)}
\begin{itemize}
\item @@$PLACEHOLDER_{59}$@@: fator de volume de formação do óleo (bbl/STB ou m$^3$/m$^3$).
\item @@$PLACEHOLDER_{60}$@@: razão gás/óleo dissolvido (scf/STB ou m$^3$/m$^3$).
\item @@$PLACEHOLDER_{61}$@@: fator de compressibilidade do gás (adimensional).
\item @@$PLACEHOLDER_{62}$@@: densidade (kg/m$^3$).
\item @@$PLACEHOLDER_{63}$@@: viscosidade (cP).
\end{itemize}

\subsection{Dicas rápidas de estudo}
\begin{itemize}
\item Verifique sempre as unidades antes de aplicar fórmulas; a maioria dos erros vem de conversões incorretas.
\item Plote curvas PVT simplificadas (@@$PLACEHOLDER_{64}$@@, @@$PLACEHOLDER_{65}$@@, @@$PLACEHOLDER_{66}$@@) para compreender tendências com pressão.
\item Use exemplos numéricos pequenos (1 STB) para fixar interpretação física dos parâmetros.
\item Ao estudar equações de estado, foque primeiro no conceito físico (equilíbrio de fases) antes da implementação numérica.
\end{itemize}


Muito obrigado(a)!
\clearpage
\subsection*{Perguntas para estudo --- Capítulo 2}
@@$PLACEHOLDER_{0}$@@

\clearpage
\subsection{Perguntas modelo — Capítulo 4}
\begin{enumerate}
\item Dissertação (Redação) — Capítulo 4.

Tema: `\texttt{Incertezas no método volumétrico: fontes, quantificação e mitigação''.

Instruções: Elabore uma redação técnica organizada em Introdução, Desenvolvimento e Conclusão. Discuta pelo menos duas fontes de incerteza e descreva um método prático para quantificá‑las (máx. 1 página).

\\vspace\{4\\baselineskip\}
\\item Questões analíticas:
\\begin\{enumerate\}[label=\\alph*)]
\\item Liste três fontes principais de incerteza no cálculo volumétrico e proponha uma forma de reduzir cada uma.
\\item Descreva em passos como montar uma simulação Monte Carlo simples para OOIP (entradas, distribuições, saídas esperadas).
\\end\{enumerate\}
\\end\{enumerate\}

\\clearpage
\\subsection\{Perguntas modelo — Capítulo 3\}
\\begin\{enumerate\}
\\item Cálculo/prático (mostre cálculos):
\\begin\{enumerate\}[label=\\alph*)]
\\item Uma amostra saturada com óleo pesa 130 g; seca pesa 105 g; @@$PLACEHOLDER_{67}$@@ g/cm@@$PLACEHOLDER_{68}$@@; @@$PLACEHOLDER_{69}$@@ cm@@$PLACEHOLDER_{70}$@@. Calcule o volume de fluido nos poros @@$PLACEHOLDER_{71}$@@ e a porosidade @@$PLACEHOLDER_{72}$@@.\\@@$PLACEHOLDER_{0}$@@

Arredondando a 3 algarismos significativos: @@$PLACEHOLDER_{73}$@@.

2) Porosidade (@@$PLACEHOLDER_{74}$@@)

Por definição:
@@$PLACEHOLDER_{1}$@@

Convertendo para percentagem e arredondando a 3 algarismos significativos: @@$PLACEHOLDER_{75}$@@.

Conclusão: a porosidade da amostra é aproximadamente @@$PLACEHOLDER_{76}$@@.

\\subsubsection\{Exemplo 2 — (Exercício 5, método de Arquímedes)\}
Dados: peso seco @@$PLACEHOLDER_{77}$@@g; peso saturado @@$PLACEHOLDER_{78}$@@g; peso aparente em água @@$PLACEHOLDER_{79}$@@g; densidade da água @@$PLACEHOLDER_{80}$@@g/cm@@$PLACEHOLDER_{81}$@@.

\\begin\{enumerate\}
\\item Volume total da amostra via empuxo (diferença entre peso em ar e peso aparente em água):
@@$PLACEHOLDER_{4}$@@
\\item Volume poroso (água/fluido preenchendo os poros):
@@$PLACEHOLDER_{5}$@@
\\item Porosidade:
@@$PLACEHOLDER_{6}$@@
\\end\{enumerate\}

Interpretação: o método de Arquímedes é útil quando o volume total não é fornecido diretamente.

\\subsection\{Exemplo rápido — compressibilidade de poros\}
Se um volume poroso @@$PLACEHOLDER_{82}$@@cm@@$PLACEHOLDER_{83}$@@ diminui @@$PLACEHOLDER_{84}$@@cm@@$PLACEHOLDER_{85}$@@ quando a pressão aumenta @@$PLACEHOLDER_{86}$@@psi, então a compressibilidade efectiva vale:
@@$PLACEHOLDER_{7}$@@
(apresente o valor em módulo como magnitude positiva quando reportar resultados).

\\subsection\{Medidas e técnicas comuns\}
\\begin\{itemize\}
\\item Análise de testemunhos (core): porosidade por pesagem, porosimetria de mercúrio, perméabilímetro de bancada.
\\item Técnicas não destrutivas: NMR, micro‑CT, imageamento 3D para distribuição de poros.
\\item Medição de pressão capilar: centrífuga e porosimetria; integrar com curvas de permeabilidade relativa.
\\item Integração: comparar core, registos de poço (neutron/densidade) e NMR para validação.
\\end\{itemize\}

\\subsection\{Glossário de símbolos (cap.3)\}
\\begin\{itemize\}
\\item @@$PLACEHOLDER_{87}$@@ — porosidade total (fração).
\\item @@$PLACEHOLDER_{88}$@@ — porosidade efectiva (fração).
\\item @@$PLACEHOLDER_{89}$@@ — volume poroso (cm@@$PLACEHOLDER_{90}$@@ ou m@@$PLACEHOLDER_{91}$@@).
\\item @@$PLACEHOLDER_{92}$@@ — volume total da amostra/elemento (cm@@$PLACEHOLDER_{93}$@@ ou m@@$PLACEHOLDER_{94}$@@).
\\item @@$PLACEHOLDER_{95}$@@ — saturações de água e óleo (frações adimensionais).
\\item @@$PLACEHOLDER_{96}$@@ — permeabilidade absoluta (Darcy, m@@$PLACEHOLDER_{97}$@@).
\\item @@$PLACEHOLDER_{98}$@@ — permeabilidade relativa (adimensional).
\\item @@$PLACEHOLDER_{99}$@@ — pressão capilar (Pa ou psi).
\\item @@$PLACEHOLDER_{100}$@@ — compressibilidade efectiva dos poros (Pa@@$PLACEHOLDER_{101}$@@ ou psi@@$PLACEHOLDER_{102}$@@).
\\end\{itemize\}

\\subsection\{Dicas rápidas de estudo\}
\\begin\{itemize\}
\\item Anote sempre as unidades e converta antes de aplicar fórmulas (campo vs SI).
\\item Modele pequenos exemplos numéricos (1 amostra; 1 STB) para fixar conceitos.
\\item Compare resultados laboratoriais com registos de poço; investigue discrepâncias.
\\item Para curvas @@$PLACEHOLDER_{103}$@@, utilize a função de Leverett @@$PLACEHOLDER_{104}$@@ para correlacionar entre rochas.
\\end\{itemize\}


Muito obrigado(a)!

\\subsection*\{Perguntas para estudo --- Capítulo 3\}
@@$PLACEHOLDER_{1}$@@
@@$PLACEHOLDER_{2}$@@

\% ---- file: $cap4_{transcription}$.md



1



ENGENHARIA DE RESERVATÓRIOS I
3º ANO – 2º SEMESTRE
2025 - 2026
1
Docente: Prof. Dr. Geraldo A. R. Ramos
1Engenharia de Reservatórios I
Geraldo Ramos, BSc, MSc, PhD
Capítulo 4 - Integração de dados e Cálculo Volumétrico de Hidrocarbonetos
\\section\{ENGENHARIA DE RESERVATÓRIOS I\}
 \\textbf\{3º Ano – 2º Semestre — 2025–2026\}

 \\textbf\{Docente:\} Prof. Dr. Geraldo A. R. Ramos

 ---

\\subsection\{Capítulo 4 — Integração de dados e Cálculo Volumétrico de Hidrocarbonetos\}

\\subsubsection\{Sumário\}
\\begin\{itemize\}
\\item 4.1 Objectivo Geral e Objectivos Específicos
\\item 4.2 Prefixos
\\item 4.3 Introdução
\\item 4.4 Definições
\\item 4.5 Métodos de cálculo de volumes originalmente existentes
\\item 4.6 Método volumétrico (parâmetros e fórmulas)
\\item 4.7 Integração de dados e incertezas
\\item 4.8 Consolidação da aula (exercícios)
\\item 4.9 Tarefas
\\item Bibliografia
\\end\{itemize\}

 ---

\\subsection\{4.1 Objetivo geral\}
Aplicar integração de dados geológicos, petrofísicos e de engenharia para estimar volumes originais de hidrocarbonetos (OOIP / OGIP) e volumes recuperáveis, considerando incertezas e apresentando resultados úteis para tomada de decisão.

\\subsubsection\{Objetivos específicos\}
\\begin\{itemize\}
\\item Definir conceitos fundamentais relacionados a volumes em reservatórios.
\\item Explicar métodos de cálculo (volumétrico, material balance, decline, probabilístico).
\\item Descrever parâmetros essenciais do método volumétrico e suas fontes de incerteza.
\\item Demonstrar integração de dados (mapas de espessura, Net‑to‑Gross, phi, Sw, Bo, Bg).
\\end\{itemize\}

\\subsection\{4.2 Prefixos e unidades (resumo)\}
\\begin\{tabular\}\{l l c\}
\\textbf\{Prefixo\} & \\textbf\{Símbolo\} & \\textbf\{Fator\} \\\\
Kilo & k & @@$PLACEHOLDER_{105}$@@ \\\\
Mega & M & @@$PLACEHOLDER_{106}$@@ \\\\
Giga & G & @@$PLACEHOLDER_{107}$@@ \\\\
Milli & m & @@$PLACEHOLDER_{108}$@@ \\\\
Micro & @@$PLACEHOLDER_{109}$@@ & @@$PLACEHOLDER_{110}$@@ \\\\
\\end\{tabular\}

\\medskip
\\textbf\{Nota:\} escolha unidades consistentes (SI ou campo — acres/ft/bbl). Ao usar fórmulas práticas, aplique os fatores de conversão corretos.

\\subsection\{4.3 Introdução\}
O método volumétrico estima o volume de fluido originalmente existente no reservatório a partir de propriedades petrofísicas e geométricas: área, espessura neta, porosidade, saturação, e fatores de volume de formação. É a primeira estimativa após descoberta e serve de base para cálculos de reservas.

\\subsection\{4.4 Definições essenciais\}
\\begin\{itemize\}
\\item @@$PLACEHOLDER_{111}$@@: volume de rocha do reservatório (m³ ou ft³);
\\item @@$PLACEHOLDER_{112}$@@: porosidade (fração);
\\item @@$PLACEHOLDER_{113}$@@: saturação de água (fração); @@$PLACEHOLDER_{114}$@@ saturação de óleo;
\\item @@$PLACEHOLDER_{115}$@@: fatores de volume de formação (reservatório → superfície);
\\item @@$PLACEHOLDER_{116}$@@: número de barris padrão (STB) ou volume de gás em condições padrão.
\\end\{itemize\}

\\subsection\{4.5 Métodos de cálculo de volumes originalmente existentes\}
\\begin\{itemize\}
\\item Método volumétrico (petrofísico/geométrico).
\\item Material balance (balanço de material) — requer dados de produção e pressão.
\\item Métodos empíricos e baseados em declínio — para reservas quando há produção.
\\item Abordagens probabilísticas (Monte Carlo) — para quantificar incerteza.
\\end\{itemize\}

\\subsection\{4.6 Método volumétrico\}
\\subsubsection\{4.6.1 Forma geral (reservoir conditions)\}
O volume de hidrocarboneto no reservatório (volume de fluido em reservatório):
@@$PLACEHOLDER_{8}$@@

Para converter para unidades de superfície usa‑se o fator de volume de formação @@$PLACEHOLDER_{117}$@@:
@@$PLACEHOLDER_{9}$@@

\\subsubsection\{4.6.2 Fórmulas práticas (exemplos)\}
\\begin\{itemize\}
\\item OOIP (stock tank barrels, usando unidades campo):
@@$PLACEHOLDER_{10}$@@
onde @@$PLACEHOLDER_{118}$@@ em acres, @@$PLACEHOLDER_{119}$@@ em ft, @@$PLACEHOLDER_{120}$@@ fração, @@$PLACEHOLDER_{121}$@@ em bbl/STB.
\\end\{itemize\}

\\begin\{itemize\}
\\item OGIP (scf ou Sm@@$PLACEHOLDER_{122}$@@ dependendo da convenção):
@@$PLACEHOLDER_{11}$@@
\\end\{itemize\}

> Observação: ao usar SI, converta @@$PLACEHOLDER_{123}$@@ (m²), @@$PLACEHOLDER_{124}$@@ (m) e fatores de conversão adequados.

\\subsubsection\{4.6.3 Estimativa de volume recuperável (reservas)\}
Reservas (recuperáveis) podem ser expressas como função das condições iniciais e finais:

Para gás (exemplo conceitual):
@@$PLACEHOLDER_{12}$@@

Para óleo (contando saturações iniciais e residuais):
@@$PLACEHOLDER_{13}$@@

Onde índices @@$PLACEHOLDER_{125}$@@ e @@$PLACEHOLDER_{126}$@@ referem‑se às condições iniciais e de abandono; @@$PLACEHOLDER_{127}$@@ e @@$PLACEHOLDER_{128}$@@ são saturação de óleo residual e fator de volume no estado residual.

\\subsubsection\{4.6.4 Reserva por fator de recuperação\}
Reservas também podem ser estimadas por:
@@$PLACEHOLDER_{14}$@@
onde @@$PLACEHOLDER_{129}$@@ é o volume originalmente em lugar e @@$PLACEHOLDER_{130}$@@ o fator de recuperação (fração).

\\subsection\{4.7 Integração de dados e incertezas\}
Passos práticos:
\\begin\{itemize\}
\\item QA/QC das fontes: mapas de batimetria, logs, cores, testes de poço, análises PVT.
\\item Construir mapas de Net‑to‑Gross, porosidade média, Sw média por setor.
\\item Calcular sensibilidade dos parâmetros-chave (A, h, φ, Sw, Bo) — análise P10/P50/P90.
\\item Usar Monte Carlo para propagar incertezas e gerar distribuições de OOIP/OGIP e reservas.
\\end\{itemize\}

Principais fontes de incerteza: interpretação de limites (cut‑offs), heterogeneidade vertical e lateral, erros de core/log, variabilidade PVT.

\\subsection\{4.8 Consolidação da aula — Exercícios\}
\\begin\{enumerate\}
\\item Dado: Área = 200 acres; $h_{net}$ = 30 ft; φ = 0.18; Swi = 0.25; Bo = 1.2 bbl/STB. Calcule OOIP (use fórmula prática).
\\item Faça uma análise de sensibilidade variando φ entre 0.15–0.22 e interprete o impacto no OOIP.
\\item Explique como Net‑to‑Gross e cut‑offs alteram a estimativa volumétrica.
\\end\{enumerate\}

\\subsection\{4.9 Tarefas\}
\\begin\{enumerate\}
\\item Reunir os dados petrofísicos e montar uma planilha de cálculo volumétrico para um caso hipotético.
\\item Rodar uma simulação Monte Carlo simples (50–100 iterações) variando φ, Sw e Bo e apresentar P10/P50/P90.
\\item Elaborar um pequeno relatório discutindo fontes de incerteza e medidas para redução de risco.
\\end\{enumerate\}

\\subsection\{Definições importantes\}
\\begin\{itemize\}
\\item \\textbf\{Net‑to‑Gross (N/G):\} fração da coluna que é considerada 'net pay' (contribui para produção).
\\item \\textbf\{Cut‑off:\} valor mínimo de porosidade/permeabilidade/saturação usado para definir zonas produtivas.
\\item \\textbf\{Fator de recuperação ($F_{R}$):\} fração do volume originalmente em lugar que se espera recuperar.
\\item \\textbf\{Área efetiva ($A_{eff}$):\} área do reservatório considerada produtiva após aplicar cut‑offs e mapas de net.
\\end\{itemize\}

\\subsection\{Exemplo prático — cálculo de OOIP (exercício 4.8)\}
Dados: Área @@$PLACEHOLDER_{131}$@@acres; @@$PLACEHOLDER_{132}$@@ft; @@$PLACEHOLDER_{133}$@@; @@$PLACEHOLDER_{134}$@@; @@$PLACEHOLDER_{135}$@@bbl/STB.

Usando a fórmula prática:
@@$PLACEHOLDER_{15}$@@
Substituindo os valores:
@@$PLACEHOLDER_{16}$@@
Calculando passo a passo:
@@$PLACEHOLDER_{17}$@@
Interpretação: aproximadamente @@$PLACEHOLDER_{136}$@@ STB (5,24 milhões de barris originalmente em lugar).

\\subsection\{Observações sobre unidades\}
\\begin\{itemize\}
\\item Fórmulas práticas usam unidades de campo (acres, ft, bbl). Para SI, converta área e espessura e aplique os fatores apropriados.
\\item Sempre verifique se @@$PLACEHOLDER_{137}$@@ e @@$PLACEHOLDER_{138}$@@ estão no mesmo sistema de unidades antes de dividir.
\\end\{itemize\}

\\subsection\{Glossário de símbolos (cap.4)\}
\\begin\{itemize\}
\\item @@$PLACEHOLDER_{139}$@@ — área (acres ou m@@$PLACEHOLDER_{140}$@@).
\\item @@$PLACEHOLDER_{141}$@@ — espessura neta (ft ou m).
\\item @@$PLACEHOLDER_{142}$@@ — porosidade (fração).
\\item @@$PLACEHOLDER_{143}$@@ — saturação de água (fração); @@$PLACEHOLDER_{144}$@@.
\\item @@$PLACEHOLDER_{145}$@@, @@$PLACEHOLDER_{146}$@@ — fatores de volume de formação do óleo e do gás.
\\item @@$PLACEHOLDER_{147}$@@, OGIP — volumes originalmente em lugar (STB para óleo; scf ou Sm@@$PLACEHOLDER_{148}$@@ para gás).
\\item @@$PLACEHOLDER_{149}$@@ — fator de recuperação (fração).
\\end\{itemize\}

\\subsection\{Dicas rápidas de estudo\}
\\begin\{itemize\}
\\item Monte uma planilha com a fórmula prática e permita variação de parâmetros (sensibilidade).
\\item Construa mapas de $A_{eff}$ e $h_{net}$ para setores diferentes e compare OOIP setorial.
\\item Use Monte Carlo para obter P10/P50/P90 e comunicar risco na estimativa.
\\end\{itemize\}


Muito obrigado(a)!

\\subsection*\{Perguntas para estudo --- Capítulo 4\}
@@$PLACEHOLDER_{3}$@@

\\clearpage
\\section\{Capítulo 6 — Simulados por Matriz\}
\\subsection\{1) Verdadeiro/Falso — Capítulo 1 (simulados)\}
@@$PLACEHOLDER_{4}$@@
\\clearpage
\\subsection\{2) Dissertação — Capítulo 4 (simulados)\}
@@$PLACEHOLDER_{5}$@@
\\clearpage
\\subsection\{3) Cálculo/prático — Capítulo 2 (simulados)\}
@@$PLACEHOLDER_{6}$@@
\\clearpage
\\subsection\{4) Cálculo/prático — Capítulo 3 (simulados)\}
@@$PLACEHOLDER_{7}$@@
\\clearpage
\\section\{Tópicos para prova\}
\\begin\{enumerate\}
\\item Verdadeiro ou falso — Capítulo 1. Para cada afirmação, assinale V (verdadeiro) ou F (falso) e justifique em até duas linhas.
\\begin\{enumerate\}[label=\\alph*)]
\\item Em reservatórios leves, a viscosidade do óleo costuma ser menor que em óleos pesados. \\hfill \\textbf\{(V/F)\}\\\\
\\textbf\{Justifique:\} \\vspace\{2\\baselineskip\}
\\item A presença de gás dissolvido no óleo sempre aumenta a mobilidade do fluido. \\hfill \\textbf\{(V/F)\}\\\\
\\textbf\{Justifique:\} \\vspace\{2\\baselineskip\}
\\item Uma drive de água (water drive) mantém a pressão do reservatório por influxo de água. \\hfill \\textbf\{(V/F)\}\\\\
\\textbf\{Justifique:\} \\vspace\{2\\baselineskip\}
\\item A solução de hidrocarbonetos pode condensar no reservatório se a pressão cair abaixo do ponto de orvalho. \\hfill \\textbf\{(V/F)\}\\\\
\\textbf\{Justifique:\} \\vspace\{2\\baselineskip\}
\\item A saturação relativa da água é igual a 1 menos a saturação do óleo. \\hfill \\textbf\{(V/F)\}\\\\
\\textbf\{Justifique:\} \\vspace\{2\\baselineskip\}
\\end\{enumerate\}

\\vspace\{4mm\}

\\item Dissertação (Redação) — Capítulo 4.

Tema: }\texttt{Incertezas no método volumétrico: fontes, quantificação e mitigação''.

Instruções: Elabore uma redação técnica organizada em Introdução, Desenvolvimento e Conclusão. Discuta pelo menos duas fontes de incerteza (ex.: variação de porosidade, estimativa de saturações, volumes de formação) e descreva um método prático para quantificá‑las (por exemplo, análise de sensibilidade ou Monte Carlo). Máximo: 1 página.

\\vspace\{6\\baselineskip\}

\\item Cálculo/prático — Capítulo 2.

Enunciado: Um ensaio PVT fornece para uma amostra padrão: volume no reservatório antes de correção @@$PLACEHOLDER_{150}$@@ bbl; volume na superfície medido @@$PLACEHOLDER_{151}$@@ STB; gás liberado @@$PLACEHOLDER_{152}$@@ scf/STB.

a) Calcule o fator de formação do óleo @@$PLACEHOLDER_{153}$@@ em bbl/STB (defina a expressão usada).  
b) Converta @@$PLACEHOLDER_{154}$@@ para m@@$PLACEHOLDER_{155}$@@/m@@$PLACEHOLDER_{156}$@@. (Use: @@$PLACEHOLDER_{157}$@@; @@$PLACEHOLDER_{158}$@@.)

Espaço para cálculos: \\vspace\{12\\baselineskip\}

\\vspace\{6mm\}

\\item Cálculo/prático — Capítulo 3.

Enunciado: Amostra de rocha saturada: massa saturada @@$PLACEHOLDER_{159}$@@; massa seca @@$PLACEHOLDER_{160}$@@; densidade do óleo @@$PLACEHOLDER_{161}$@@; volume total da amostra @@$PLACEHOLDER_{162}$@@.

a) Calcule o volume de fluido nos poros @@$PLACEHOLDER_{163}$@@ e a porosidade @@$PLACEHOLDER_{164}$@@ (\\\%).  
b) Apresente os passos e respostas com 3 algarismos significativos.

Espaço para cálculos: \\vspace\{12\\baselineskip\}

\\end\{enumerate\}
\\clearpage
\\section\{Banco de Questões — Matriz completa\}
Este banco reúne perguntas extraídas e adaptadas dos capítulos 1–4 para revisão e geração de provas.

\\subsection\{Capítulo 1 — Verdadeiro/Falso e questões curtas\}
\\begin\{enumerate\}
\\item Em reservatórios leves, a viscosidade do óleo costuma ser menor que em óleos pesados. (V/F)
\\item A presença de gás dissolvido no óleo sempre aumenta a mobilidade do fluido. (V/F)
\\item Um mecanismo de water drive mantém a pressão por influxo de água. (V/F)
\\item A saturação relativa da água é igual a 1 menos a saturação do óleo. (V/F)
\\item Explique em 3 linhas o papel de uma armadilha geológica na acumulação de hidrocarbonetos.
\\item Descreva a diferença entre Sistema de Produção e Sistema Petrolífero (3 linhas).
\\end\{enumerate\}

\\subsection\{Capítulo 2 — Cálculo/prático\}
\\begin\{enumerate\}
\\item (PVT) Dado: @@$PLACEHOLDER_{165}$@@ bbl; @@$PLACEHOLDER_{166}$@@ STB; @@$PLACEHOLDER_{167}$@@ scf/STB. Calcule @@$PLACEHOLDER_{168}$@@ e converta @@$PLACEHOLDER_{169}$@@ para m@@$PLACEHOLDER_{170}$@@/m@@$PLACEHOLDER_{171}$@@.
\\item (Z‑factor) Quando é aceitável tratar o gás como ideal? Explique como obter @@$PLACEHOLDER_{172}$@@ na prática.
\\item (Interpretação) Dado um conjunto de curvas @@$PLACEHOLDER_{173}$@@, discuta o que uma queda acentuada de @@$PLACEHOLDER_{174}$@@ indica sobre o sistema.
\\end\{enumerate\}

\\subsection\{Capítulo 3 — Cálculo/prático\}
\\begin\{enumerate\}
\\item Calcule porosidade a partir de massas e densidade do fluido (ex.: 130 g/105 g, @@$PLACEHOLDER_{175}$@@ g/cm@@$PLACEHOLDER_{176}$@@, @@$PLACEHOLDER_{177}$@@ cm@@$PLACEHOLDER_{178}$@@).
\\item Método de Arquímedes: calcule @@$PLACEHOLDER_{179}$@@, @@$PLACEHOLDER_{180}$@@ e @@$PLACEHOLDER_{181}$@@ a partir dos dados de peso em água e peso saturado.
\\item Calcule a porosidade média para um conjunto de amostras e discuta a variabilidade.
\\end\{enumerate\}

\\subsection\{Capítulo 4 — Dissertação e análise\}
\\begin\{enumerate\}
\\item Redija uma dissertação técnica sobre }\texttt{Incertezas no método volumétrico'' (máx. 1 página).
\\item Liste e explique três medidas práticas para reduzir incerteza na estimativa de OOIP.
\\item Esboce os passos para uma análise de sensibilidade (P10/P50/P90) usando Monte Carlo.
\\end\{enumerate\}

\\vspace\{6mm\}

\\clearpage
\\section\{Exercícios e Tarefas — Reunidos (arquivos da pasta)\}
\\begin\{comment\}
\\subsection\{$prova_{banco_questoes_toda_matriz}$.md\}
@@$PLACEHOLDER_{0}$@@

\\subsection\{$prova_{banco_questoes_cap1}$-4.md\}
@@$PLACEHOLDER_{1}$@@

\\subsection\{$provas_{vf_cap1}$.md\}
@@$PLACEHOLDER_{2}$@@

\\subsection\{$provas_{calc_cap2}$.md\}
@@$PLACEHOLDER_{3}$@@

\\subsection\{$provas_{calc_cap3}$.md\}
@@$PLACEHOLDER_{4}$@@

\\subsection\{$provas_{dissertacao_cap4}$.md\}
@@$PLACEHOLDER_{5}$@@

\\subsection\{$exercicios_{resolvidos}$.md\}
@@$PLACEHOLDER_{6}$@@
\\end\{comment\}

\\clearpage
\\section*\{Notas\}
Os blocos acima reúnem os exercícios e tarefas extraídos dos ficheiros Markdown presentes na pasta }Estudar/matéria/revisão para a prova\texttt{. Alguns itens podem duplicar conteúdos já incluídos nas secções de consolidação dos capítulos; mantive cada fonte como subseção identificada para controlo e revisão.

Se preferir, posso:
\\begin\{itemize\}
\\item Converter estes blocos em LaTeX nativo (enumerates/itemize e math) em vez de verbatim,
\\item Deduplicar questões repetidas automaticamente,
\\item Inserir apenas exercícios novos que não estejam já no compêndio,
\\item Gerar e compilar o PDF atualizado imediatamente.
\\end\{itemize\}
 
\% Conteúdo do usuário integrado nas seções correspondentes.

\\clearpage

\\textbf\{Solução (passo a passo)\}

1) Volume de fluido nos poros ($V_{f}$)

Usa‑se a relação entre massa de fluido e densidade:

@@$PLACEHOLDER_{2}$@@

Substituindo os valores:

@@$PLACEHOLDER_{3}$@@

Arredondando para 3 algarismos significativos: $V_{f}$ ≈ 29.8 cm³.

2) Porosidade (φ)

Por definição: @@$PLACEHOLDER_{182}$@@.

Substituindo:

@@$PLACEHOLDER_{4}$@@

Convertendo para percentagem e arredondando a 3 algarismos significativos: φ ≈ 16.5 \%.


\\textbf\{Resultados finais\}

- Volume de fluido nos poros: $V_{f}$ ≈ 29.8 cm³
- Porosidade: φ ≈ 16.5 \%

\\textbf\{Observações\}
- Unidades usadas: g, cm³, g/cm³. Para SI, converta massa (kg) e volume (m³) consistentemente.  
- Método alternativo (método de Arquímedes) e outros exemplos estão disponíveis em }$exercicios_{resolvidos}$.md\texttt{ e no material do Cap.3.

# Capítulo 1 — Afirmações V/F (prontas)

Marque V ou F e justifique brevemente.

1. Reservatórios leves produzem óleo com baixa densidade e alta mobilidade. (V/F)

2. A presença de gás dissolvido no óleo reduz a viscosidade e aumenta a mobilidade. (V/F)

3. Um water drive mantém a pressão do reservatório por influxo de água do aquífero. (V/F)

4. A expansão do gás dissolvido no óleo contribui para a manutenção da pressão no solution gas drive. (V/F)

5. A existência de um gas cap (camada de gás livre) ajuda a manter a pressão por expansão do gás. (V/F)

6. O ponto de bolha indica a pressão em que começa a aparecer gás livre no óleo. (V/F)

7. Net‑to‑Gross não afeta a estimativa de OOIP quando a porosidade é conhecida. (V/F)

8. A saturação irreductível de água (Swi) reduz o volume de óleo disponível. (V/F)

9. A viscosidade do óleo não influencia a mobilidade do fluido. (V/F)

10. As fórmulas práticas de OOIP devem ser aplicadas somente após verificar unidades e fatores de conversão. (V/F)


Justifique cada resposta em uma ou duas linhas abaixo da afirmação.

# Resumo por Capítulo — Engenharia de Reservatórios I (VERSÃO ULTRA‑DETALHADA)

Este ficheiro contém uma compilação exaustiva de fórmulas, definições, derivações e exemplos numéricos para os capítuos 1–5. Use como referência técnica aprofundada durante a preparação.

Índice
- Capítulo 1 — Sistema petrolífero e conceitos fundamentais
- Capítulo 2 — Propriedades dos fluidos (PVT): equações, correlações e cálculos
- Capítulo 3 — Propriedades de rochas: porosidade, permeabilidade, capilaridade, testes
- Capítulo 4 — Cálculo volumétrico (OOIP / OGIP): fórmulas de campo e SI, sensibilidade, Monte Carlo
- Capítulo 5 — Equação de Balanço de Materiais (EBM): formulação, linearização, p/z e exemplos
- Anexos: constantes, fatores de conversão, valores típicos, checklist de passos práticos


\\textbf\{Observação sobre unidades\}
- Indique sempre o sistema: CAMPO (acres, ft, bbl, scf, psi) ou SI (m, m³, Pa). Mantenha consistência.
- Fatores de conversão essenciais estão no Anexo.


## Capítulo 1 — Sistema petrolífero e conceitos fundamentais

1. Definições essenciais
- Sistema petrolífero: rocha geradora + rocha selante + armadilha + migração + sincronismo. 
- Sistema de produção: conjunto de poços, tubulações, elevação artificial, separadores, tratamento e escoamento.

2. Relações e equações úteis
- Equilíbrio hidrostático (coluna vertical):
@@$PLACEHOLDER_{18}$@@
(usar densidade local em cada camada quando heterogênea)

... (conteúdo adicional do compêndio omitido por brevidade nesta inserção verbatim)


# Compressibilidade (fator) Z — Cheatsheet extremo

Objetivo: referência completa e passo‑a‑passo para entender, calcular e aplicar o fator de compressibilidade @@$PLACEHOLDER_{183}$@@ (gás real), com métodos práticos (Standing–Katz, EoS), derivadas úteis e um exemplo numérico com código Python.


\\textbf\{Resumo rápido\}
- \\textbf\{Definição:\} @@$PLACEHOLDER_{184}$@@, mede o desvio do gás em relação ao gás ideal (@@$PLACEHOLDER_{185}$@@ ideal).
- \\textbf\{Usos:\} converter entre densidade e pressão/temperatura reais; calcular o fator de volume @@$PLACEHOLDER_{186}$@@; avaliar compressibilidade isotérmica do gás; PVT e balanço de massas.


\\textbf\{8) Código Python (PR EoS) — cálculo de Z e compressibilidade por diferença finita\}

Instalar dependência mínima: }pip install numpy\texttt{

}`\texttt{python
import numpy as np

R = 8.314462618

def $Z_{PR}$(p, T, Tc, Pc, omega):
	# p in Pa, T in K, Tc in K, Pc in Pa
	a = 0.45724 \\emph\{ R\\textbf\{2 \} Tc\}2 / Pc
	b = 0.07780 \\emph\{ R \} Tc / Pc
	kappa = 0.37464 + 1.54226\\emph\{omega - 0.26992\}omega**2
	alpha = (1 + kappa*(1 - np.sqrt(T/Tc)))**2
	A = a \\emph\{ alpha \} p / (R\\textbf\{2 * T\}2)
	B = b \\emph\{ p / (R \} T)
	# Coeffs cubic: Z^3 - (1-B) Z^2 + (A - 3B^2 - 2B) Z - (A B - B^2 - B^3) = 0
	coeffs = [1.0, -(1 - B), (A - 3\\emph\{B\}B - 2\\emph\{B), -(A\}B - B\\emph\{B - B\}B*B)]
	roots = np.roots(coeffs)
	$real_{roots}$ = np.real(roots[np.isclose(roots.imag, 0, atol=1e-8)])
	if len($real_{roots}$) == 0:
		raise RuntimeError('No real root found for Z')
	# gas root = largest real root
	Z = np.max($real_{roots}$)
	return float(Z)

def $gas_{compressibility}$(p, T, Tc, Pc, omega, dp=1e5):
	Z1 = $Z_{PR}$(p, T, Tc, Pc, omega)
	Z2 = $Z_{PR}$(p + dp, T, Tc, Pc, omega)
	dZdp = (Z2 - Z1) / dp
	$c_{g}$ = 1.0/p - (1.0/Z1)*dZdp
	return Z1, dZdp, $c_{g}$

if __name__ == '__main__':
	# Example: methane at p=5 MPa, T=350 K
	Tc = 190.564
	Pc = 4.5992e6
	omega = 0.011
	p = 5e6
	T = 350.0
	Z, dZdp, cg = $gas_{compressibility}$(p, T, Tc, Pc, omega, dp=1e5)
	print('Z =', Z)
	print('dZ/dp =', dZdp, 'Pa^-1')
	print('$c_{g}$ =', cg, 'Pa^-1 (', cg*1e6, 'MPa^-1 )')
}`\texttt{

Saída esperada aproximada (exemplo que usamos manualmente): }Z ≈ 0.948\texttt{, }$c_{g}$ ≈ 2.19e-7 Pa^-1\texttt{.


FIM — ficheiro gerado como versão ULTRA‑DETALHADA. Se desejar, posso:
- (A) sobrepor }$resumo_{capitulos}$.md` com esta versão;
- (B) gerar PDF/LaTeX desta versão; ou
- (C) extrair cartões Anki (Q/A) automaticamente a partir das definições e fórmulas.

Arquivo criado automaticamente pelo assistente.
\end{Verbatim}

\end{document}
